% How a term is stored and how the terms are split (deck, the four words used in the design part;
% README 2.2 and 3.2), with a three-qubit example.
%   1  a term is a coefficient and a Pauli string; the string is kept as two bit vectors (x | z), one bit
%      pair per qubit: I = (0,0), X = (1,0), Z = (0,1), Y = (1,1). X I Z is 100 | 001.
%   2  a row r is the parity of the bits that a mask selects; the mask 000 | 111 selects the three z bits
%      (0, 0, 1) of X I Z, so r = 1
%   3  a sector is the set of terms whose rows have the same values (two rows give sectors 00, 01, 10, 11)
%   4  a window is a run of consecutive gates that keep these rows unchanged
% The coefficient 0.31 and the string X I Z are an illustration; the figure carries no measured numbers.
\begin{tikzpicture}[x=1cm, y=1cm, font=\rmfamily\footnotesize]
  \tikzset{bb/.style={draw=black!70, line width=0.35pt, anchor=north west, inner sep=0pt, minimum width=0.55cm,
                      minimum height=0.54cm, text height=0.3cm, text depth=0.09cm, font=\rmfamily\small},
           hh/.style={draw=black!70, line width=0.35pt, anchor=north west, inner sep=0pt, fill=black!10,
                      minimum height=0.36cm, text height=0.2cm, text depth=0.07cm, font=\rmfamily\scriptsize},
           lg/.style={draw=black!70, line width=0.35pt, anchor=north west, inner sep=0pt, minimum width=0.6cm,
                      minimum height=0.42cm, text height=0.26cm, text depth=0.08cm, font=\rmfamily\footnotesize},
           tb/.style={draw=black!70, line width=0.3pt, anchor=west, inner sep=0pt, minimum width=0.3cm, minimum height=0.3cm},
           gt/.style={draw=black!70, line width=0.35pt, anchor=west, inner sep=0pt, minimum width=0.46cm, minimum height=0.46cm}}
  \node[layer, anchor=north west, minimum width=13.9cm, minimum height=7.2cm] at (0,0) {};
  % ================= 1: a term and its bits =================
  \node[num] at (0.47,-0.5) {1};
  \node[cttl, anchor=west] at (0.68,-0.5) {Term};
  \node[hh, minimum width=1.3cm] at (0.35,-1.0) {Coefficient};
  \node[hh, minimum width=1.65cm] at (1.65,-1.0) {Pauli String};
  \node[bb, fill=white, minimum width=1.3cm] at (0.35,-1.36) {0.31};
  \foreach \i/\p in {0/X, 1/I, 2/Z}
    \node[bb, fill=white] at ({1.65+0.55*\i},-1.36) {\p};
  \draw[arr] (3.45,-1.63) -- (4.27,-1.63);
  \node[hh, minimum width=1.65cm] at (4.4,-1.0) {x};
  \node[hh, minimum width=1.65cm] at (6.3,-1.0) {z};
  \foreach \i/\b in {0/1, 1/0, 2/0}
    \node[bb, fill=white] at ({4.4+0.55*\i},-1.36) {\b};
  \foreach \i/\b in {0/0, 1/0, 2/1}
    \node[bb, fill=rowc!35] at ({6.3+0.55*\i},-1.36) {\b};
  % the two bits of one qubit
  \node[lg, fill=black!10, minimum width=0.95cm] at (10.15,-0.8) {Pauli};
  \node[lg, fill=black!10, minimum width=0.95cm] at (10.15,-1.22) {x};
  \node[lg, fill=black!10, minimum width=0.95cm] at (10.15,-1.64) {z};
  \foreach \i/\p/\x/\z in {0/I/0/0, 1/X/1/0, 2/Z/0/1, 3/Y/1/1}
  {
    \node[lg, fill=black!10] at ({11.1+0.6*\i},-0.8) {\p};
    \node[lg, fill=white] at ({11.1+0.6*\i},-1.22) {\x};
    \node[lg, fill=white] at ({11.1+0.6*\i},-1.64) {\z};
  }
  % ================= 2: a row =================
  \node[num] at (0.47,-2.69) {2};
  \node[cttl, anchor=west] at (0.68,-2.69) {Row $r$};
  \node[note, anchor=east] at (4.3,-2.69) {Mask};
  \foreach \i in {0,1,2}
  {
    \node[bb, fill=white] at ({4.4+0.55*\i},-2.42) {0};
    \node[bb, fill=rowc] at ({6.3+0.55*\i},-2.42) {1};
  }
  \draw[arr] (8.1,-2.69) -- (8.87,-2.69);
  \node[key, anchor=west] at (8.95,-2.69) {$r$ = Parity of the Selected Bits = 1};
  \draw[black!55, dashed, line width=0.5pt] (0.2,-3.55) -- (13.7,-3.55);
  % ================= 3: sectors =================
  \node[num] at (0.47,-4.1) {3};
  \node[cttl, anchor=west] at (0.68,-4.1) {Sector};
  \foreach \k/\id/\f/\n in {0/00/sa/3, 1/01/sb/2, 2/10/sc/3, 3/11/sd/2}
  {
    \node[draw=black!70, line width=0.45pt, fill=\f!45, anchor=north west, minimum width=1.5cm, minimum height=1.55cm] at ({0.35+1.6*\k},-4.5) {};
    \node[ittl, anchor=north] at ({1.1+1.6*\k},-4.53) {Sector \id};
    \foreach \i in {1,...,\n}
      \node[tb, fill=\f] at ({1.1+1.6*\k-0.15*\n+0.3*(\i-1)},-5.5) {};
  }
  \node[note, anchor=base] at (3.5,-6.62) {Terms Whose Rows Have the Same Values};
  % ================= 4: a window =================
  \node[num] at (7.57,-4.1) {4};
  \node[cttl, anchor=west] at (7.78,-4.1) {Window};
  \node[draw=navyf, line width=0.7pt, fill=white, anchor=west, minimum width=3.36cm, minimum height=0.76cm] at (7.45,-5.275) {};
  \foreach \k in {0,...,5}
    \node[gt, fill=genc] at ({7.55+0.54*\k},-5.275) {};
  \foreach \k in {0,...,3}
    \node[gt, fill=genc!40] at ({10.97+0.54*\k},-5.275) {};
  \node[font=\rmfamily\footnotesize] at (13.4,-5.275) {$\cdots$};
  \node[note, anchor=base] at (10.45,-6.62) {Consecutive Gates That Keep These Rows Unchanged};
\end{tikzpicture}
